\section{Classes R and M3}\label{sec:reverse}

This section excludes the two classes whose obstruction is local: an entirely reverse period
(class $\cls{R}$) and a mixed period with a maximal reverse run of length other than two (class
$\cls{M3}$). Both exclusions hold for every normal vertex period, whatever its cost and its
number of pieces. They rest on two facts about a few consecutive pieces. Across three reverse
switches the active margin of an interior piece cannot be a minimum of the four
(\cref{lem:reverse-four}). Across two switches that return to the same label, the active margins
are strictly log-concave (\cref{lem:three-arc}). In a run of reverse switches these two facts
leave room for exactly one configuration, a run of two switches around a strict minimum
(\cref{prop:run-structure}). That configuration is the input of the remaining classes. The
section ends with the proof of ellipticity (\cref{prop:ellipticity}), which uses it.

\subsection{The reverse-four principle}\label{sec:rev-four}

Consider four consecutive pieces $k,k+1,k+2,k+3$ with three reverse switches between them, so
that their labels are $i,\,i-1,\,i+1,\,i$. Name the margins at the switch $\tau_{k+1}$ in the label
order $(i,i+1,i-1)$ as $(P,Q,Z)$, and the factors of the two interior pieces as $b=q_{k+1}$ and
$c=q_{k+2}$. Then $P=a_k$ and $Z=a_{k+1}$. Along piece $k+1$ (label $i-1$) the label $i$ grows and
$i+1$ decreases, so at $\tau_{k+2}$
\[
  A:=m_{k+2}(i+1)=\frac{Z+Q}{b}-Z=a_{k+2},\qquad B:=m_{k+2}(i)=(Z+P)\,b-Z,
\]
and along piece $k+2$ (label $i+1$) the label $i$ decreases, so $a_{k+3}=(A+B)/c-A$.

Suppose the third piece is a weak minimum: $a_{k+2}\le a_k,a_{k+1},a_{k+3}$. In these variables
this reads
\begin{equation}\label{eq:rev-min}
  A\le Z,\qquad bA+(b-1)Z\le B,\qquad (2c-1)A\le B.
\end{equation}
The obstruction is a variation of the period (\cite{Notes},
\path{GLOBAL_REVERSE_PERIOD_EXCLUSION}): insert short pieces of the two
inactive controls at the two interior switches (the control $i+1$ at $\tau_{k+1}$ and the
control $i$ at $\tau_{k+2}$), and change the logarithms of the factors of the four pieces at
rates chosen so that the end state and the end frame do not move. To first order the normalized
cost $\cost/\sqrt{12}$ then changes by $\frac{Z}{4}N(A,B,Z,b,c)$ times the size of the variation, where
\begin{multline}\label{eq:rev-N}
  N(A,B,Z,b,c)=b(b-1)A\bigl(Ac\,\eta(c)-B\,\zeta(c)\bigr)\\
  +(c-1)\bigl(bA(B+bZ)\,\eta(b)-B(bA+Z)\,\zeta(b)\bigr),
\end{multline}
with $\eta(t)=t-1-\log t$ and $\zeta(t)=t\log t-t+1$, both positive for $t>1$. Under the
minimum hypothesis~\eqref{eq:rev-min} this derivative is negative
(\cref{lem:rev-N}), so the variation lowers the cost, while the maximum principle says that the
inserted controls are not better than the active ones and the variation cannot lower the cost to
first order. We do not construct the variation. Its content is an exact linear identity between
switching functions at the switches (\cref{lem:rev-identities}), and the maximum principle enters
only through the signs of these switching functions.

\begin{lemma}[Sign of the cost derivative]\label{lem:rev-N}
If $b,c>1$, $A,Z>0$ and~\eqref{eq:rev-min} holds, then $N(A,B,Z,b,c)<0$.
\end{lemma}

The proof (\cref{app:rev-N}) rests on the inequality $(2t+1)\,\eta(t)<(t+2)\,\zeta(t)$ for $t>1$,
which is the Pad\'e bound $\log t>3(t^2-1)/(t^2+4t+1)$.

\begin{lemma}[The reverse-four identities]\label{lem:rev-identities}
In the situation above, with $\chi_j(k)$ the switching functions at the switch $\tau_k$
(\cref{def:costate-coordinates}),
\begin{align}
  B(c-1)\,\chi_{i+1}(k+1)+bD\,\chi_i(k+2)&=-\sqrt3\,Z\,N(A,B,Z,b,c),\label{eq:rev-third}\\
  (b-1)B\,\chi_{i-1}(k+3)+cD\,\chi_i(k+2)&=-\sqrt3\,A\,N(Z,B,A,c,b),\label{eq:rev-second}
\end{align}
where $D=A(c-1)+Z(b-1)$.
\end{lemma}

\begin{proof}
Express the costate coordinates at $\tau_{k+2}$ and $\tau_{k+3}$ through those at $\tau_{k+1}$ by
the transport of \cref{lem:costate-transport}(b) over the two interior pieces. The left sides
are then linear in the six coordinates $\alpha_j(k+1),\beta_j(k+1)$, with coefficients rational
in $P,Q,Z,b,c$ and linear in $\log b$, $\log c$. Five linear relations hold among these
coordinates: the switching functions of the active labels vanish at
$\tau_{k+1}$ (labels $i$ and $i-1$), at $\tau_{k+2}$ (label $i+1$) and at $\tau_{k+3}$ (label $i$),
and $\ip{\Lambda_R}{X}=0$ at $\tau_{k+1}$. Each identity is a combination of these five relations
with polynomial coefficients, listed in \cref{app:rev-identities}; after clearing denominators
both sides are the same polynomial in the variables and in $\log b$, $\log c$.
\end{proof}

The coefficients $B(c-1)$ and $bD$ in~\eqref{eq:rev-third} are the weights of the two
inserted pieces in the variation described above, and $N$ is its cost derivative.

\begin{lemma}[Reverse-four principle]\label{lem:reverse-four}
Let four consecutive pieces $k,\dots,k+3$ of a normal vertex period have three reverse switches
between them. Then
\begin{enumerate}[label=\textup{(\alph*)}]
\item it is not the case that $a_{k+2}\le a_k$, $a_{k+2}\le a_{k+1}$ and $a_{k+2}\le a_{k+3}$;
\item it is not the case that $a_{k+1}\le a_k$, $a_{k+1}\le a_{k+2}$ and $a_{k+1}\le a_{k+3}$.
\end{enumerate}
\end{lemma}

\begin{proof}
(a) The labels $i+1$ at $\tau_{k+1}$ and $i$ at $\tau_{k+2}$ are inactive on both adjacent pieces,
so their switching functions are $\le0$ (\cref{lem:costate-transport}(c)). The weights
$B(c-1)$ and $bD$ are positive, so the left side of~\eqref{eq:rev-third} is $\le0$
and $N(A,B,Z,b,c)\ge0$. The weak minimum at the third piece is~\eqref{eq:rev-min}, and
\cref{lem:rev-N} gives $N<0$.

(b) A weak minimum at the second piece reads $Z\le A$, $cZ+(c-1)A\le B$ and $(2b-1)Z\le B$: these
are the conditions~\eqref{eq:rev-min} with $(A,Z,b,c)$ replaced by $(Z,A,c,b)$, so
$N(Z,B,A,c,b)<0$. The switching functions $\chi_{i-1}(k+3)$ and $\chi_i(k+2)$ are $\le0$, and
\eqref{eq:rev-second} gives the contradiction as before.
\end{proof}

Part (b) also follows from (a) by a reflection of the trajectory combined with a reversal of
time, which reverses the order of the four pieces and keeps the signs of the switches
(\cite{Notes}, \path{MIXED_GLOBAL_MINIMUM_STRUCTURE}, \S\S1--4). The second identity
\eqref{eq:rev-second} avoids this transformation of trajectories.

\subsection{Class R}\label{sec:rev-R}

\begin{theorem}[Class $\cls{R}$]\label{thm:class-R}
No normal vertex period has every switch reverse.
\end{theorem}

\begin{proof}
The active margins are $N$-periodic, so one of them, $a_m$, is a global minimum. Every switch is
reverse, so the four pieces $m-2,\dots,m+1$ have three reverse switches between them, and
$a_m\le a_{m-2},a_{m-1},a_{m+1}$. This contradicts \cref{lem:reverse-four}(a) for $k=m-2$.
\end{proof}

The window may run through the seam, and it may contain the same piece of the period twice when
$N<4$; the twisted indexing of \cref{sec:periods} makes no distinction. The argument uses neither
the cost nor the number of pieces, and \cref{sec:terminal} uses it also for the entirely reverse
words with five pieces.

\subsection{Three pieces with alternating labels}\label{sec:rev-three}

When the switch after piece $k$ is followed by a switch back to the label of piece $k$, the
middle piece is a turn, and its two outer pieces have the same control. The switching function
of that control vanishes at both ends of the middle piece and is $\le0$ in between. This forces
a relation between the three active margins.

\begin{lemma}[Three-arc log-concavity and turn ordering]\label{lem:three-arc}
Let $k,k+1,k+2$ be consecutive pieces of a normal vertex period.
\begin{enumerate}[label=\textup{(\alph*)}]
\item If $i_{k+2}=i_k$, then $a_ka_{k+2}<a_{k+1}^2$; in particular $a_{k+1}$ is not at most both
  $a_k$ and $a_{k+2}$.
\item If the switch $k$ is reverse and the switch $k+1$ forward, then $a_k<a_{k+2}$; if the switch
  $k$ is forward and the switch $k+1$ reverse, then $a_{k+2}<a_k$.
\end{enumerate}
\end{lemma}

\begin{proof}
(b) After a reverse switch the label $i_k$ is the growing label of piece $k+1$, so
$a_{k+2}=(a_{k+1}+a_k)\,q_{k+1}-a_{k+1}>a_k$ by \cref{lem:one-arc}; after a forward switch it is
the decreasing label and $a_{k+2}=(a_{k+1}+a_k)/q_{k+1}-a_{k+1}<a_k$.

(a) Write $a=a_{k+1}$ and $q=q_{k+1}$. Suppose first that the switch $k$ is reverse, so that the
outer label $i=i_k$ is the growing label of the middle piece, with margin $b=a_k$ at its start, and
$a_{k+2}=(a+b)q-a$. Since
\[
  a^2-b\bigl((a+b)q-a\bigr)=(a+b)(a-bq),
\]
the claim is $bq<a$. By \cref{lem:costate-transport} and \cref{app:rev-one-arc}, the switching
function of the label $i$ at the intermediate factor $s\in[1,q]$ is
\begin{gather*}
  \chi_i(s)=2s\Bigl(Y(s-1)+\frac{\sqrt3a}{2}F(s)\Bigr),\\
  F(s)=(a+b)\log s-(a+b)(s-1)+a\frac{(s-1)^2}{s},
\end{gather*}
for a number $Y$. It vanishes at $s=1$ and $s=q$ and is $\le0$ in between. The vanishing at $q$
fixes $Y$, and $\chi_i\le0$ just before $q$ makes the derivative at $q$ nonnegative. Together
these say $(q-1)F'(q)\ge F(q)$, that is,
\[
  (q-1)(2aq-a+bq)\ge(a+b)\,q^2\log q.
\]
If $bq\ge a$, this fails: the right side minus the left side equals
$a\bigl(q^2\log q-(q-1)(2q-1)\bigr)+b\bigl(q^2\log q-q(q-1)\bigr)$; the second bracket is
positive because $q\log q>q-1$; if the first bracket is negative, replace $a$ by the larger
$bq$, which leaves $bq^2\bigl((q+1)\log q-2(q-1)\bigr)>0$. So $bq<a$.

If the switch $k$ is forward, the outer label is the decreasing label of the middle piece, with
margin $c=a_k$, and $a_{k+2}=(a+c)/q-a$. Now $a^2-c\bigl((a+c)/q-a\bigr)=(a+c)(a-c/q)$, and the
claim is $c<aq$. Here $\chi_i(s)=-(2/s)\bigl(Y(1-s^{-1})+\frac{\sqrt3a}{2}h(s)\bigr)$ with
$h(s)=(a+c)\log s-a(s-1)-c(1-s^{-1})$, and $h(1)=h'(1)=0$. The sign at the start of the piece
gives $Y\ge0$. If $aq\le c$, then
\[
  h(q)=(c-aq)\bigl(\log q-(1-q^{-1})\bigr)+a\bigl((q+1)\log q-2(q-1)\bigr)>0,
\]
so $\chi_i(q)<0$, contradicting $\chi_i(q)=0$.
\end{proof}

The proof uses one end of the middle piece in each case: the right end when the outer label grows,
the left end when it decreases. \Cref{fig:reverse-estimates}(a) shows the switching function of
the growing case for an admissible factor and for one with $bq>a$. Another proof of (a) works on
the edge of the simplex that contains the two controls: in a clock in which the state moves at
unit speed, one transforms the costate accordingly and compares the two switching abscissas
(\cite{Notes}, \path{MIXED_GLOBAL_MINIMUM_STRUCTURE}~(12)).

\subsection{Reverse runs}\label{sec:rev-runs}

\begin{proposition}[Run structure]\label{prop:run-structure}
Let $P$ be a mixed normal vertex period, and let the switches $k,\dots,k+L-1$ form a maximal
reverse run. Then $L=2$, and the active margin of the centre piece $k+1$ is strictly smaller
than those of the four pieces around it:
\[
  a_{k+1}<a_{k-1},\qquad a_{k+1}<a_k,\qquad a_{k+1}<a_{k+2},\qquad a_{k+1}<a_{k+3}.
\]
\end{proposition}

\begin{proof}
The run is preceded by the forward switch $k-1$ and followed by the forward switch $k+L$.

At the entry, the switches $k-1,k$ are forward then reverse, so $a_{k+1}<a_{k-1}$ by
\cref{lem:three-arc}(b), and $i_{k+1}=i_{k-1}$. If $a_k\le a_{k+1}$, then $a_k$ would be at most
both $a_{k-1}$ and $a_{k+1}$, against \cref{lem:three-arc}(a). So $a_{k+1}<a_k$.

At the exit, the switches $k+L-1,k+L$ are reverse then forward, so $a_{k+L-1}<a_{k+L+1}$, and
\cref{lem:three-arc}(a) on the pieces $k+L-1,k+L,k+L+1$ gives $a_{k+L-1}<a_{k+L}$.

If $L=1$, the entry gives $a_{k+1}<a_k$ and the exit $a_k<a_{k+1}$. If $L\ge3$, let $m$ be a
piece of $k,\dots,k+L$ with the smallest active margin. It is not $k$, since $a_{k+1}<a_k$, nor
$k+L$, since $a_{k+L-1}<a_{k+L}$. If $m=k+1$, the pieces $k,\dots,k+3$ have three reverse switches
and the second one is a weak minimum, against \cref{lem:reverse-four}(b). If
$k+2\le m\le k+L-1$, the pieces $m-2,\dots,m+1$ have three reverse switches and the third one is a
weak minimum, against \cref{lem:reverse-four}(a). So $L=2$. The entry and the exit give the first
three inequalities, and \cref{lem:three-arc}(b) at the switches $k+1,k+2$ gives
$a_{k+1}<a_{k+3}$.
\end{proof}

The proof does not use mixedness beyond the existence of the run: a maximal reverse run is
bounded by forward switches by definition, so a period that has one is mixed.

\begin{theorem}[Class $\cls{M3}$]\label{thm:class-M3}
No normal vertex period is mixed with a maximal reverse run of length other than two.
\end{theorem}

\begin{proof}
This is the first conclusion of \cref{prop:run-structure}.
\end{proof}

In a mixed period every reverse switch therefore belongs to a pattern $+\,-\,-\,+$ of switch
signs, whose middle piece is a centre, and every centre is a strict local minimum of the active
margins. The other classes use these facts through the following corollary.

\begin{corollary}[Mixed periods]\label{cor:rev-mixed}
Let $P$ be a mixed normal vertex period.
\begin{enumerate}[label=\textup{(\alph*)}]
\item Every reverse switch $k$ is the first or the second switch of a pattern of signs
  $+\,-\,-\,+$: either the switches $k-1,k,k+1,k+2$ or the switches $k-2,k-1,k,k+1$ have these
  signs.
\item At a centre $k$, $a_k<a_{k-1}$ and $a_k<a_{k+1}$.
\item Every switch norm satisfies $d_k\ge3/32$.
\end{enumerate}
\end{corollary}

\begin{proof}
(a) and (b) restate \cref{prop:run-structure}. For (c), at a forward switch
$d_k\ge j_k/2\ge3/32$ by \cref{lem:switch-norms}(a). By (a), a reverse switch is the outgoing
switch of a turn (the piece $k$ after the forward switch $k-1$) or the incoming switch of a turn
(the piece $k+1$ before the forward switch $k+1$), and \cref{lem:switch-norms}(b) gives
$d_k\ge3/32$ in both cases.
\end{proof}

\subsection{Proof of ellipticity}\label{sec:rev-ellipticity}

We prove \cref{prop:ellipticity}: a mixed normal vertex period has $\gamma^2>0$. Suppose
$\gamma^2\le0$. The idea is to read the switch vectors $y_n$ of \cref{lem:per-switch-vectors}
projectively, through a coordinate that changes monotonically from one switch to the next.
Because $y_{N+n}=\pm y_n$, the coordinate returns to its value after one period; but at every
centre it strictly increases, and nowhere does it decrease.

\begin{lemma}[Three Gram vectors]\label{lem:rev-three-gram}
Let $S$ be a symmetric $2\times2$ matrix with $\det S\le0$, and let $y,y',y''$ be vectors with
$y^{\mathsf T}Sy$, $y'^{\mathsf T}Sy'$, $y''^{\mathsf T}Sy''$ all positive. Then
$(y^{\mathsf T}Sy')(y^{\mathsf T}Sy'')(y'^{\mathsf T}Sy'')>0$.
\end{lemma}

\begin{proof}
For two vectors, $(y^{\mathsf T}Sy)(y'^{\mathsf T}Sy')-(y^{\mathsf T}Sy')^2=\det S\cdot\det(y,y')^2\le0$.
The Gram determinant of three vectors of the plane vanishes:
$d\,d'\,d''-d\,c^2-d'\,b^2-d''a^2+2abc=0$ in the obvious notation. With the three two-vector
inequalities, $2abc\ge d\,d'\,d''>0$.
\end{proof}

\begin{lemma}[Signs of the cross pairings]\label{lem:rev-cross-signs}
Suppose $\gamma^2\le0$. At a forward same-sign piece $k$, $e_k>0$. At a centre $k$ with
$a_k<a_{k-1}$, $a_k<a_{k+1}$ and $d_{k-1},d_k\ge3/32$, $e_k<0$.
\end{lemma}

\begin{proof}
At a same-sign piece, \cref{lem:stationarity} expresses $e_k$ through $d_{k-1}$, $d_k$ and the
active margins. At a forward same-sign piece both switches are forward, and the lower bounds
$d\ge j/2$ of \cref{lem:switch-norms}(a) give $e_k\ge X_k$ with
\[
  X_k=\frac1{16}\Bigl(-1+2a\Bigl(\frac1{a+p}+\frac1{a+n}\Bigr)-\frac{a^2}{a^2+4}
  +\log\frac{(a+p)(a+n)}{a^2+4}\Bigr),
\]
$a=a_k$, $p=a_{k-1}$, $n=a_{k+1}$. Here the logarithm is $\log q_k>0$ by
\cref{lem:per-factors}, and $a^2/(a^2+4)<2a/(a+p)$ because $ap<4$; so $X_k>-1/16$. On the other
hand $\gamma^2\le0$ and \cref{lem:gram} give $e_k^2\ge d_{k-1}d_k\ge(3/32)^2$, and $3/32>1/16$, so
$e_k\ge3/32>0$. At a centre the switches are reverse, the upper bounds $d\le j/2$ give
$e_k\le X_k$, and now the logarithm is $-\log q_k<0$ and $2a/(a+p)+2a/(a+n)<2$ by the strict
minimum, so $X_k<1/16$ and $e_k\le-3/32$.
\end{proof}

\begin{proof}[Proof of \cref{prop:ellipticity}]
Suppose $\gamma^2\le0$. By \cref{cor:rev-mixed}(c) every $d_n=y_n^{\mathsf T}Sy_n$ is at least
$3/32$. Let $h(v)=y_0^{\mathsf T}Sv$, a nonzero linear form, and let $l$ be the linear form with
$h(u)l(v)-h(v)l(u)=c\det(u,v)$ for a constant $c>0$. By \cref{lem:rev-three-gram} with
$y=y_0$, $h(y_n)h(y_m)\,(y_n^{\mathsf T}Sy_m)>0$ for all $n,m$; in particular $h(y_n)\ne0$, and the
projective coordinate $z_n=l(y_n)/h(y_n)$ is defined, with
\[
  z_{n+1}-z_n=\frac{c\,\det(y_n,y_{n+1})}{h(y_n)\,h(y_{n+1})}.
\]
By \cref{lem:per-switch-vectors}, at a turn $y_{n+1}=-y_n$ and the increment is $0$. At a forward
same-sign piece $\det(y_n,y_{n+1})=1/K_{n+1}>0$ and $y_n^{\mathsf T}Sy_{n+1}=e_{n+1}>0$
(\cref{lem:rev-cross-signs}), so $h(y_n)h(y_{n+1})>0$ and the increment is positive. At a centre
$\det(y_n,y_{n+1})=-1/K_{n+1}<0$ and $e_{n+1}<0$, using the strict minimum and the bounds
$d\ge3/32$ of \cref{cor:rev-mixed}; so $h(y_n)h(y_{n+1})<0$ and the increment is again positive.
A mixed period has a reverse switch, hence a centre (\cref{cor:rev-mixed}(a)), so the increments
over one period sum to a positive number. But $y_N=\pm y_0$ and $z$ depends only on the line of
the vector, so the sum is $z_N-z_0=0$.
\end{proof}

The same telescoping can be organized with a factorization of the indefinite form $S$ in place
of \cref{lem:rev-three-gram} (\cite{Notes}, \path{MIXED_GLOBAL_COSTATE_ELLIPTIC}, \S\S5--6). For an entirely
forward period every piece is a forward same-sign piece and the same telescoping gives
$\gamma^2>0$.
